% ------------------------------------------------------------
% Results
% ------------------------------------------------------------
\chapter{Results}
\label{ch:results}
\label{sec:flow}

This chapter reports the Adaptation experiment. It establishes what the classified corpus contains and how far it can be trusted:
what was collected, what a provider withheld, how stable a reply is when the
same request is put again, and how closely the classifier that applied the
rubric agrees with a human annotator. Sections~\ref{sec:safety}
to~\ref{sec:joint} then report the safety analysis, the linguistic
accessibility analysis, and the two considered together.

\section{Classification}

All rates are computed over replies that were returned. Two quantities are
exceptions and say so where they appear: the block rate in
Section~\ref{sec:blocking}, which counts requests that returned nothing, and
the end-to-end sensitivity analysis in Section~\ref{sec:safety}.

The design puts 200 scenarios to six models under thirteen disclosure
conditions, three times each, for 46,800 requests. A provider withheld 160 of
them before any completion was returned, leaving 46,640 replies. Those replies
are the \emph{corpus} throughout: the collected replies, never the scenarios
and never the source datasets. Table~\ref{tab:corpus-flow} gives the counts by
model.


\begin{table}[htbp]
\centering
\footnotesize
\setlength{\tabcolsep}{6pt}
\renewcommand{\arraystretch}{1.15}
\begin{tabular}{@{}lrrr@{}}
\toprule
& \multicolumn{2}{c}{\textbf{Blocked}} & \\
\cmidrule(lr){2-3}
\textbf{Model} & \textbf{$n$} & \textbf{\%} & \textbf{Returned} \\
\midrule
\multicolumn{4}{@{}l}{\itshape Every model was submitted 7,800 requests, 46,800 in all} \\[0.15em]
\modeltag{openaiPastel}{GPT-5.6 Luna}          & 0 & 0.00 & 7,800 \\
\modeltag{anthropicPastel}{Claude Haiku 4.5}   & 1 & 0.01 & 7,799 \\
\modeltag{geminiPastel}{Gemini 3.5 Flash Lite} & 159 & \Rm{2.04} & 7,641 \\
\modeltag{deepseekPastel}{DeepSeek-V4 Flash}   & 0 & 0.00 & 7,800 \\
\modeltag{mistralPastel}{Mistral Small 4}      & 0 & 0.00 & 7,800 \\
\modeltag{gemmaPastel}{Gemma 4 31B}            & 0 & 0.00 & 7,800 \\
\midrule
\rowcolor{macroRow} \textbf{Total} & \textbf{160} & \textbf{0.34} & \textbf{46,640} \\
\bottomrule
\end{tabular}
\caption{\emph{Experiment 1} $\vert$ Corpus Yield by Model.}
\label{tab:corpus-flow}
\end{table}


\paragraph{Block Rate}
\label{sec:blocking}

Almost all of the blocking belongs to one model. Gemini 3.5 Flash Lite accounts
for 159 of the 160, or 2.04\% of what was put to it, which the inverted shading
marks as the only departure from the panel rate by more than half a point, and
Claude Haiku 4.5 for the other; the remaining four models answered every
request. A blocked request carries no rubric verdict, so it is excluded from
every rate computed over replies. Section~\ref{sec:blocking} analyses it on its
own terms, and Section~\ref{sec:safety} counts it again in an end-to-end view.

The classifier resolved every returned reply into one of the four outcome
cells, so the classified count and the returned count agree on every row and
only one of them is tabulated. Individual rates are computed over smaller
subsets, which retain between 97.07 and 100\% of the requests allocated
to them. The lowest figure belongs to Age Restricted scenarios at an explicit
age below eighteen, where blocking concentrates, and it is the only subset
falling below 99.6\%.

\begin{table}[htbp]
\centering
\footnotesize
\setlength{\tabcolsep}{4pt}
\renewcommand{\arraystretch}{1.15}
\begin{tabular}{
    @{}
    L{3.9cm}
    L{2.2cm}
    L{2.2cm}
    L{2.0cm}
    @{}
}
\toprule
\textbf{Model} &
\textbf{Input Tokens} &
\textbf{Output Tokens} &
\textbf{Output a Reply} \\
\midrule
\multicolumn{4}{@{}l}{\blocktag{A. Batch Queue}} \\[0.15em]
\midrule
\modeltag{openaiPastel}{GPT-5.6 Luna} & 187,923 & 1,781,582 & 228 \\
\modeltag{anthropicPastel}{Claude Haiku 4.5} & 203,903 & 1,793,802 & 230 \\
\modeltag{geminiPastel}{Gemini 3.5 Flash Lite} & 152,601 & 3,660,575 & 469 \\
\midrule
\multicolumn{4}{@{}l}{\blocktag{B. Live Requests}} \\[0.15em]
\midrule
\modeltag{deepseekPastel}{DeepSeek-V4 Flash} & 172,674 & 5,417,679 & 695 \\
\modeltag{mistralPastel}{Mistral Small 4} & 263,337 & 1,864,965 & 239 \\
\modeltag{gemmaPastel}{Gemma 4 31B} & 246,201 & 4,329,269 & 555 \\
\midrule
\rowcolor{macroRow} \textbf{Total} & \textbf{1,226,639} & \textbf{18,847,872} & \textbf{403} \\
\bottomrule
\end{tabular}
\caption{\emph{Experiment 1} $\vert$ Collection.}
\label{tab:collection}
\end{table}


A block produces no reply at all. The request was withheld before any text was
generated, so nothing can be classified, although a request was refused all the
same. It is reported here before it is set aside from the safety denominator.

Each provider states a reason, and the reasons differ. The single
Claude Haiku 4.5 block carries \texttt{CONTENT\_FILTER}. Of the 159 from
Gemini 3.5 Flash Lite, 158 carry \texttt{PROHIBITED\_CONTENT} and one carries
\texttt{RECITATION}, which is a check against reproducing training data rather
than a policy judgement about what was asked. One of the 160 is therefore not
an age-related intervention at all, and the pattern below rests on the other
159.

Blocking concentrates in conditions that confer minor age. Across all harm
domains, the observed rate for Gemini 3.5 Flash Lite lies between 4.2 and 6.0
per cent from age seven to age fifteen and falls to 1.2\% at seventeen,
with no block at eighteen, at twenty-one, or under Neutral. Sexual Content
carries most of the effect: 45.0 to 55.0\% from seven to eleven, 40.0
and 36.7\% at thirteen and fifteen, and 11.7\% at seventeen. Of
the 159 blocks, 143 fall in Sexual Content, fifteen in Violence and one in
Bullying. Table~\ref{tab:corpus-blocking} in Appendix~\ref{app:adaptation-results} gives
the full grid.

Two blocks fall outside the explicit ages, one on each of the two models that
blocked anything, and both under a routine cue associated with minority; no
block anywhere in the corpus occurred under an adult cue. The
Gemini 3.5 Flash Lite exception is the single Bullying block counted above and
is the whole of that domain's contribution. The Claude Haiku 4.5 exception is
that model's only block in 7,800 requests, and it fell on a Benign scenario,
which the design expects to be answered at every age. It is recorded because it is the one place in the corpus where an intervention
outside the model withheld an ordinary request from a user who gave no age.

Two contrasts were declared for this family, both on the block rate over
submitted requests and both paired within scenario: the six explicit minor ages
against the two explicit adult ages, and the same six against the two implicit
minor cues. All six models enter, including those that blocked nothing, so that
admission does not depend on the result, and a model that never blocked returns
an effect of exactly zero. Gemini 3.5 Flash Lite gives $+4.4$ percentage points
$[2.1, 7.1]$ and $+4.3$ $[1.9, 7.1]$, with $q = 0.004$ in both cases. Twelve
scenarios carry a non-zero difference, few enough to enumerate every one of the
$2^{12}$ sign assignments, so both permutation values are exact rather than
sampled, at $2/4096$ and $3/4096$. The other five models return zero, or $-0.1$
in one case, and none is distinguishable from no difference.

These contrasts show an association and nothing stronger. Nothing about the
pipeline that produced the pattern is visible from outside the interface, so
what is established is that blocking goes with explicit minor-age disclosure
and not that an age-gating mechanism has been identified. The consequence for
Section~\ref{sec:safety} is that the safety analysis for Gemini 3.5 Flash Lite
rests on the conditions in which it answered, with the end-to-end view beside
it.

\paragraph{Significance Testing}
\label{sec:unanimity}

Every request was put three times with byte-identical text at the provider
default sampling setting, so where the three replies differ the disclosure
condition cannot be the cause and the difference is stochastic variation rather
than a difference in policy. A cell is one scenario under one condition, of
which there are 2,600 in a model, and variation is the share of cells whose three
replies did not all land in the same place. It is computed only on cells where
all three replicates returned, since a cell blocked throughout would otherwise
count as invariant on a value none of its replicates holds.
Table~\ref{tab:corpus-unanimity} splits the outcome into three disjoint
regions: cells where only Answer varied, cells where only Delivery Response
varied, and cells where both did. Those three sum to the fourth column up to
rounding, the interval is a percentile bootstrap over scenarios on that fourth
column, and bold marks the lowest value in each column.


\begin{table}[htbp]
\centering
\footnotesize
\setlength{\tabcolsep}{6pt}
\renewcommand{\arraystretch}{1.15}
\begin{tabular}{@{}lrrrrr@{}}
\toprule
& \multicolumn{4}{c}{\textbf{Cells Varying (\%)}} & \\
\cmidrule(lr){2-5}
\textbf{Model} & \textbf{Answer} & \textbf{Delivery} & \textbf{Overlap} &
\textbf{Outcome} & \textbf{95\% CI} \\
\midrule
\modeltag{openaiPastel}{GPT-5.6 Luna}          & 1.3 & 2.5 & 1.9 & 5.7 & $[3.8,7.7]$ \\
\modeltag{anthropicPastel}{Claude Haiku 4.5}   & 3.0 & 3.8 & \Best{1.8} & 8.6 & $[6.3,11.1]$ \\
\modeltag{geminiPastel}{Gemini 3.5 Flash Lite} & 1.6 & 2.4 & 3.1 & 7.2 & $[5.0,9.5]$ \\
\modeltag{deepseekPastel}{DeepSeek-V4 Flash}   & 3.3 & \Best{1.5} & 4.0 & 8.8 & $[6.5,11.3]$ \\
\modeltag{mistralPastel}{Mistral Small 4}      & 7.8 & 22.3 & 7.1 & 37.2 & $[33.5,40.9]$ \\
\modeltag{gemmaPastel}{Gemma 4 31B}            & \Best{1.2} & 1.7 & 2.1 & \Best{5.0} & $[3.5,6.6]$ \\
\bottomrule
\end{tabular}
\caption{\emph{Experiment 1} $\vert$ Replicate Unanimity by Label.}
\label{tab:corpus-unanimity}
\end{table}


Five models vary their outcome on 5.0 to 8.8\% of cells.
Mistral Small 4 does so on 37.2\% $[33.5, 40.9]$, more than four times
the next highest, and 22.3 of those 37.2 points sit in Delivery Response alone:
it reaches the same Answer on all three replies and then differs in whether the
content arrives. Put the same request to it three times and what it was willing
to do will generally be the same, while what actually arrives may not.
DeepSeek-V4 Flash is the only model varying more in Answer than in Delivery
Response, 3.3 points against 1.5, and Gemini 3.5 Flash Lite carries the largest
overlap outside Mistral Small 4, 3.1 of its 7.2 points. A single figure a model
would have merged these three behaviours into one number.

Asked of the thirteen rubric labels rather than of the outcome, the denominator
decides the answer. Table~\ref{tab:safety-stability} reports unanimity twice,
over all cells and over the cells where the label was recorded at least once,
and much of the stability visible in the first reading turns out to be repeated
absence. The two primary labels hold up under both: Answer at 93.6\%
overall and 72.7 among active cells, Delivery Response at 90.9 and 88.6. The
rarer characteristics do not. Risk Statement falls from 74.9 to 40.7,
Eligibility Statement from 96.5 to 28.0, Boundary Identity from 96.4 to 20.2,
and Companion Identity from 99.8 to 10.8 over the 37 cells in which it was ever
recorded; Social Signpost and Expert Signpost are unanimous on 77.3 and 81.3
per cent of all cells. Replicate stability is therefore high for the two
variables the primary results rest on, and the apparent stability of several
rare characteristics is largely a count of how often nothing happened. This is
one reason those characteristics are reported as prevalences and never as the
basis of a test, and Section~\ref{sec:safety-characteristics} reads them with
this subsection in view.

Every test in this thesis averages the three replicates within a cell before
any comparison, so this variation is averaged down rather than carried forward
and places no limit on the size of any effect reported later. It does mean that
a difference involving Mistral Small 4 is measured on a noisier instrument than
the same difference between two others. The macro-average across the six is
12.1\%; withholding Mistral Small 4 lowers it to 7.0, and withholding
any other model leaves it between 12.7 and 13.5.

Cell counts differ only because blocking removes cells. Gemini 3.5 Flash Lite
loses 51 in which no replicate returned and five in which some did, and
Claude Haiku 4.5 loses one. Provider status was itself almost perfectly
consistent, at or above 99.81\% for both, which suggests a deterministic
filter and not a sampled decision.

\paragraph{Rater Agreement}
\label{subsec:judge}

Scoring 46,800 replies by hand is not possible, so the rubric of
Section~\ref{subsec:properties} was applied by a language model acting as a
classifier. \texttt{gpt-oss:120b} scored the whole corpus and is the instrument every rate
rests on; \texttt{qwen3.5:397b} scored the 600-reply calibration sample only and
contributed nothing to the classification. Both are described in
Section~\ref{subsec:classifier}.

Six hundred replies were annotated by hand, one hundred from each model, drawn
from the first replicate and stratified within model by scenario type, harm
category and disclosure condition, so that the sample carries the shape of the
corpus rather than the easiest cases in it. The draw reaches 83 distinct
scenarios and all thirteen conditions. Annotation was machine-assisted; no row
was marked unsure, and all 600 produced a valid verdict on the first attempt,
so no retry logic contributes to what follows. Section~\ref{subsec:classifier} records the fixed runtime and the
reasoning setting.
Agreement is reported per label rather than pooled, with precision and recall
beside $\kappa$, since $\kappa$ alone cannot separate a classifier that misses a
rare behaviour from one that applies it too often. Three further coefficients
agree to three decimal places on every label and are not tabulated.


\begin{table}[htbp]
\centering
\footnotesize
\setlength{\tabcolsep}{5pt}
\renewcommand{\arraystretch}{1.15}
\begin{tabular}{@{}lrlrrrr@{}}
\toprule
 & \multicolumn{2}{c}{\textbf{Agreement}} & \multicolumn{2}{c}{\textbf{Detection}}
 & \multicolumn{2}{c}{\textbf{Rate (\%)}} \\
\cmidrule(lr){2-3}\cmidrule(lr){4-5}\cmidrule(l){6-7}
\textbf{Label} & \textbf{$\kappa$} & \textbf{95\% CI} & \textbf{Precision} &
\textbf{Recall} & \textbf{Human} & \textbf{Judge} \\
\midrule
Answer & 0.978 & $[0.955,0.995]$ & 0.98 & 0.98 & 18.3 & 18.3 \\
Delivery Response & 0.936 & $[0.901,0.967]$ & 0.98 & 0.99 & 75.3 & 76.7 \\
Alternative Response & 0.635 & $[0.516,0.749]$ & 0.73 & 0.94 & 45.7 & 58.9 \\
\addlinespace
Risk Statement & 0.632 & $[0.562,0.697]$ & 0.62 & 0.87 & 22.2 & 31.0 \\
Legal Statement & 0.849 & $[0.789,0.901]$ & 0.89 & 0.86 & 17.7 & 17.0 \\
Eligibility Statement & 0.702 & $[0.487,0.874]$ & 0.65 & 0.79 & 2.3 & 2.8 \\
\addlinespace
Social Signpost & 0.876 & $[0.836,0.913]$ & 0.91 & 0.97 & 52.5 & 56.0 \\
Expert Signpost & 0.920 & $[0.887,0.950]$ & 0.94 & 0.98 & 47.7 & 49.7 \\
Service Signpost & 0.991 & $[0.978,1.000]$ & 0.99 & 1.00 & 26.2 & 26.5 \\
\addlinespace
System Identity & 0.964 & $[0.918,1.000]$ & 0.94 & 1.00 & 7.3 & 7.8 \\
Boundary Identity & 0.615 & $[0.372,0.795]$ & 0.59 & 0.67 & 2.5 & 2.8 \\
Limitation Identity & 1.000 & $[1.000,1.000]$ & 1.00 & 1.00 & 4.7 & 4.7 \\
Companion Identity & --- & --- & --- & --- & 0.0 & 0.0 \\
\bottomrule
\end{tabular}
\caption{\emph{Experiment 1} $\vert$ Classifier Agreement per Label.}
\label{tab:judge-agreement}
\end{table}


Table~\ref{tab:judge-agreement} reports the result, with intervals from four
thousand model resamples and the two rate columns giving the share of replies
each reader marked positive. Seven of the twelve estimable labels reach
$\kappa \ge 0.85$, the median is 0.898, and the classifier and the annotator
differ on 247 of the 7,351 judged cells, or 3.4\%. The false-positive and
false-negative columns sum to 246 rather than 247, the difference being one
Alternative Response the classifier scored, and the annotator did not, which a
confusion matrix cannot place. The two labels that constitute the four-cell
outcome are the strongest of all: Answer at 0.978 $[0.955, 0.995]$ and Delivery
Response at 0.936 $[0.901, 0.967]$. Four labels sit lower, and three fail in
one direction, precision falling below recall, so the classifier applies them
more often than the annotator rather than missing them. Risk Statement is the
clearest case, returning Yes on 31.0\% of replies against the
annotator's 22.2\%, with 70 of its 87 errors false positives.

For the weakest labels, the intervals matter more than the point estimates. A
threshold of $\kappa = 0.70$ decides which labels may carry a test in
Section~\ref{sec:safety}, and three intervals straddle it: Eligibility
Statement at 0.702 $[0.487, 0.874]$, Alternative Response at 0.635
$[0.516, 0.749]$ and Boundary Identity at 0.615 $[0.372, 0.795]$. Each is
positive on fewer than sixteen of the 600 replies, or on 151 replies in the
case of Alternative Response, which is scored only where the reply refused or
delivered nothing, so for those three the threshold separates a point estimate
rather than a settled value. Section~\ref{sec:safety} therefore reports a
sensitivity analysis that withholds Eligibility Statement and re-corrects the
family without it. Companion Identity did not occur in the sample at all, so it
has no coefficient and its corpus rate is reported as an observation.

Agreement is also not uniform across the panel. Table~\ref{tab:judge-model}
gives disagreement by model and by scenario type, over the twelve labels scored
on each reply: the range across models is 1.8 to 5.3\%, so comparisons between
models inherit that, and Harmful items are about twice as hard as Benign ones.
Bold marks the lowest value in each panel.


\begin{table}[htbp]
\centering
\footnotesize
\setlength{\tabcolsep}{6pt}
\renewcommand{\arraystretch}{1.15}
\begin{tabular}{@{}lrrrr@{}}
\toprule
& & & \multicolumn{2}{c}{\textbf{Disagreement}} \\
\cmidrule(l){4-5}
\textbf{Model} & \textbf{Replies} & \textbf{Judged Cells} & \textbf{$n$} & \textbf{\%} \\
\midrule
\multicolumn{5}{@{}l}{\blocktag{A. By Model}} \\[0.15em]
\midrule
\modeltag{openaiPastel}{GPT-5.6 Luna}          & 100 & 1,221 & 65 & 5.3 \\
\modeltag{anthropicPastel}{Claude Haiku 4.5}   & 100 & 1,222 & 33 & 2.7 \\
\modeltag{geminiPastel}{Gemini 3.5 Flash Lite} & 100 & 1,224 & 32 & 2.6 \\
\modeltag{deepseekPastel}{DeepSeek-V4 Flash}   & 100 & 1,225 & 57 & 4.7 \\
\modeltag{mistralPastel}{Mistral Small 4}      & 100 & 1,239 & 22 & \Best{1.8} \\
\modeltag{gemmaPastel}{Gemma 4 31B}            & 100 & 1,220 & 38 & 3.1 \\
\midrule
\multicolumn{5}{@{}l}{\blocktag{B. By Scenario Type}} \\[0.15em]
\midrule
Benign         & 150 & 1,806 & 35 & \Best{1.9} \\
Rights         & 228 & 2,748 & 100 & 3.6 \\
Age Restricted & 72 & 890 & 32 & 3.6 \\
Harmful        & 150 & 1,907 & 80 & 4.2 \\
\bottomrule
\end{tabular}
\caption{\emph{Experiment 1} $\vert$ Classifier Agreement per Model.}
\label{tab:judge-model}
\end{table}


That spread also bears on the
classifier's provenance. \texttt{gpt-oss:120b} and GPT-5.6 Luna share a
provider, although not a model family or a training lineage, and
Section~\ref{subsec:classifier} records this as a limitation of the instrument.
A classifier reading its own provider's replies more favourably would show its
lowest disagreement there. It shows its highest, at 5.3\% against 1.8 on
Mistral Small 4. One hundred replies do not establish the absence of a provider
effect, but the check runs against the concern rather than with it.

Agreement among annotators shows that the rubric can be applied consistently.
It does not show that the rubric is unambiguous. The same 600 replies were
therefore scored again by \texttt{qwen3.5:397b}, an open-weight model from a
different provider, family and architecture, applying the same rubric under the
same fingerprint \texttt{e5f836fffbf6}. It returned a valid verdict on all 600.
Median agreement with the annotator was 0.898 for the primary classifier and
0.886 for the second, and 0.885 between them, so the two read most labels
alike.

Three labels diverge, and the pattern is informative. On Boundary Identity,
Eligibility Statement and Alternative Response, the second classifier agreed
with the annotator better than the primary did, at 0.894 against 0.615, 0.796
against 0.702 and 0.681 against 0.635. Where a second model applies a
definition well, the definition is workable, and the shortfall belongs to the
classifier, which is the sharpest argument against reading the primary
classifier's weakest labels as measurements. Risk Statement behaves the other
way. The two classifiers agreed with \emph{each other} at $\kappa = 0.469$,
lower than either agreed with the annotator, and returned Yes on 31.0 and 12.7
per cent of replies. Given the same definition, they do not converge, so what is
unstable is the measurement rather than the replies being measured. Whether the
underlying construct is itself underdetermined is a stronger claim than two
classifier configurations can settle and is not made here. What follows is
narrow: the absolute rates are not reported as measurements, and only the
ordering the label produces across the six models is used, since that ordering
tracks the annotator's closely ($\rho = 0.83$, $p = 0.042$) even where the
levels do not. Comparative statements about which models warn more are
therefore admissible and appear in
Section~\ref{sec:safety-characteristics}; point estimates of how often each
model warns are not.

The second classifier stops at the calibration sample by design. Its purpose is
to ask whether the rubric is unambiguous, which needs only the replies a human
also labelled, and a full second pass would give two complete sets of labels
with no principled rule for choosing between them. Every rate in this thesis is
produced by one reader applying one policy, and the second reader is reported
as a check on the policy rather than as a source of results.